\section{Presheaves of algebraic structures}
\label{section-presheaves-structures}

\noindent
Let us clarify the definition
of presheaves of algebraic structures.
Suppose that $\mathcal{C}$ is a category and
that $F : \mathcal{C} \to \textit{Sets}$ is
a faithful functor. Typically $F$ is a ``forgetful''
functor. For an object $M \in \Ob(\mathcal{C})$
we often call $F(M)$ the {\it underlying set} of the
object $M$. If $M \to M'$ is a morphism in $\mathcal{C}$
we call $F(M) \to F(M')$ the {\it underlying map of sets}.
In fact, we will often not distinguish between an object
and its underlying set, and similarly for morphisms.
So we will say a map of sets $F(M) \to F(M')$
is a {\it morphism of algebraic structures}, if it is
equal to $F(f)$ for some morphism $f : M \to M'$
in $\mathcal{C}$.

\medskip\noindent
In analogy with Definition \ref{definition-abelian-presheaves}
above a ``presheaf of objects of $\mathcal{C}$'' could be
defined by the following data:
\begin{enumerate}
\item a presheaf of sets $\mathcal{F}$, and
\item for every open $U \subset X$ a choice
of an object $A(U) \in \Ob(\mathcal{C})$
\end{enumerate}
subject to the following conditions (using the phraseology above)
\begin{enumerate}
\item for every open $U \subset X$ the set $\mathcal{F}(U)$
is the underlying set of $A(U)$, and
\item for every $V \subset U \subset X$ open
the map of sets $\rho_V^U: \mathcal{F}(U) \to \mathcal{F}(V)$
is a morphism of algebraic structures.
\end{enumerate}
In other words, for every $V \subset U$ open in $X$
the restriction mappings $\rho^U_V$ is the image
$F(\alpha^U_V)$ for some unique morphism
$\alpha^U_V : A(U) \to A(V)$ in the category $\mathcal{C}$.
The uniqueness is forced by the condition that $F$ is
faithful; it also implies that
$\alpha^U_W = \alpha^V_W \circ \alpha^U_V$
whenever $W \subset V \subset U$ are open in $X$.
The system $(A(-), \alpha^U_V)$ is what we will define as a
presheaf with values in $\mathcal{C}$ on $X$, compare
Sites, Definition \ref{sites-definition-presheaf}.
We recover our presheaf of sets $(\mathcal{F}, \rho_V^U)$
via the rules $\mathcal{F}(U) = F(A(U))$ and
$\rho_V^U = F(\alpha_V^U)$.

\begin{definition}
\label{definition-presheaf-values-in-category}
Let $X$ be a topological space.
Let $\mathcal{C}$ be a category.
\begin{enumerate}
\item A {\it presheaf $\mathcal{F}$ on $X$ with values in $\mathcal{C}$}
is given by a rule which assigns to every open $U \subset X$
an object $\mathcal{F}(U)$ of $\mathcal{C}$
and to each inclusion $V \subset U$
a morphism $\rho_V^U : \mathcal{F}(U) \to \mathcal{F}(V)$
in $\mathcal{C}$ such that whenever $W \subset V \subset U$
we have $\rho_W^U = \rho_W^V \circ \rho_V^U$.
\item A {\it morphism $\varphi : \mathcal{F} \to \mathcal{G}$
of presheaves with value in $\mathcal{C}$} is given by a
morphism $\varphi : \mathcal{F}(U) \to \mathcal{G}(U)$
in $\mathcal{C}$ compatible with restriction morphisms.
\end{enumerate}
\end{definition}

\begin{definition}
\label{definition-underlying-presheaf-sets}
Let $X$ be a topological space. Let $\mathcal{C}$ be a category.
Let $F : \mathcal{C} \to \textit{Sets}$ be a faithful functor.
Let $\mathcal{F}$ be a presheaf on $X$ with values in $\mathcal{C}$.
The presheaf of sets $U \mapsto F(\mathcal{F}(U))$
is called the {\it underlying presheaf of sets of $\mathcal{F}$}.
\end{definition}

\noindent
It is customary to use the same letter $\mathcal{F}$ to denote
the underlying presheaf of sets, and this makes
sense according to our discussion preceding
Definition \ref{definition-presheaf-values-in-category}.
In particular, the phrase ``let $s \in \mathcal{F}(U)$''
or ``let $s$ be a section of $\mathcal{F}$ over $U$'' signifies
that $s \in F(\mathcal{F}(U))$.

\medskip\noindent
This notation and these definitions apply in particular to:
{\it Presheaves of (not necessarily abelian) groups, rings, modules
over a fixed ring, vector spaces over a fixed field, } etc and
{\it morphisms between these}.

